\chapter{Ciała uporządkowane i zupełność}
\label{ma1:ch:ordered-fields}

W poprzednim rozdziale skonstruowaliśmy liczby rzeczywiste dwiema metodami. Teraz oddzielimy własności wynikające z samych aksjomatów ciała uporządkowanego od tych, które wymagają zupełności. W szczególności wykażemy, że zupełność wyznacza ciało liczb rzeczywistych jednoznacznie z dokładnością do izomorfizmu, oraz skonstruujemy pierwiastki bez korzystania z ciągłości.

\section{Ciała uporządkowane}
\begin{definicja}[Ciało uporządkowane]
\label{ma1:def:ordered-field}
Ciałem uporządkowanym nazywamy ciało $F$ wyposażone w porządek liniowy $\le$, spełniający dla wszystkich $x,y,z\in F$ warunki
\begin{enumerate}[label=(\roman*)]
\item $x\le y\Longrightarrow x+z\le y+z$;
\item $0\le x,\ 0\le y\Longrightarrow 0\le xy$.
\end{enumerate}
\end{definicja}
\begin{lemat}[Reguły rachunku nierówności]
\label{ma1:lem:ordered-rules}
Niech $F$ będzie ciałem uporządkowanym. Wówczas $1>0$, kwadrat każdego niezerowego elementu jest dodatni, a dla $x<y$ i $z>0$ zachodzi $xz<yz$. Dla $z<0$ zachodzi $xz>yz$. Jeżeli $0<x<y$, to $0<1/y<1/x$.
\end{lemat}
\begin{proof}
Dla $x\ne0$ albo $x>0$, albo $-x>0$. W obu przypadkach $x^2=(-x)^2\ge0$ i $x^2\ne0$, stąd $x^2>0$. W szczególności $1=1^2>0$. Jeśli $x<y$, to $y-x>0$ na mocy zgodności porządku z dodawaniem (równość po dodaniu oznaczałaby $x=y$). Dla $z>0$ iloczyn $(y-x)z$ jest nieujemny i niezerowy, zatem $yz-xz>0$. Dla $z<0$ stosujemy wynik do $-z>0$ i zmieniamy znaki.
Dla $x>0$ odwrotność $1/x$ nie może być ujemna, gdyż wtedy iloczyn $x(1/x)=1$ byłby ujemny; nie może też być zerem. Jeżeli $0<x<y$, mnożymy nierówność $x<y$ przez dodatni element $1/(xy)$, otrzymując $1/y<1/x$.
\end{proof}
\begin{twierdzenie}[Stożek dodatni]
\label{ma1:thm:positive-cone}
Niech $F$ będzie ciałem. Porządki czyniące $F$ ciałem uporządkowanym odpowiadają bijektywnie podzbiorom $P\subseteq F$ takim, że $P+P\subseteq P$, $PP\subseteq P$ oraz $F$ jest rozłączną sumą zbiorów $P$, $\{0\}$ i $-P$. Odpowiednio $P=\{x:x>0\}$, a $x<y$ wtedy i tylko wtedy, gdy $y-x\in P$.
\end{twierdzenie}
\begin{proof}
Dla danego porządku domknięcie $P$ na dodawanie wynika z $0<x<x+y$ dla $x,y>0$, a na mnożenie z lematu \ref{ma1:lem:ordered-rules}. Rozkład rozłączny wynika z liniowości porządku.
Odwrotnie, dla danego $P$ określamy $x\le y$ przez $y-x\in P\cup\{0\}$. Zwrotność jest bezpośrednia. Jeśli $x\le y$ i $y\le x$, to albo $x=y$, albo zarówno $y-x$, jak i jego przeciwność należą do $P$, wbrew rozłączności. Przechodniość wynika z $(z-y)+(y-x)=z-x$ i domknięcia $P\cup\{0\}$ na dodawanie. Rozkład rozłączny daje porównywalność dowolnych dwóch elementów. Dodanie tego samego elementu nie zmienia różnicy, a iloczyn elementów nieujemnych jest nieujemny z domknięcia $P$ na mnożenie. Obie konstrukcje są wzajemnie odwrotne.
\end{proof}
Własności wartości bezwzględnej w dowolnym ciele uporządkowanym udowodniono w sekcji \ref{ma1:sec:rational-cauchy}. Będziemy używać ich bez powtarzania dowodów.
\begin{przyklad}
Ciała $\Q$ i $\R$ z poprzednich rozdziałów są uporządkowane. Natomiast ciało $\mathbb F_2=\{0,1\}$ z działaniami modulo $2$ nie dopuszcza takiego porządku: z $1>0$ wynikałoby $1+1>0$, lecz $1+1=0$.
\end{przyklad}

\section{Kanoniczne liczby wymierne}
\begin{twierdzenie}
\label{ma1:thm:rational-embedding}
Dla każdego ciała uporządkowanego $F$ istnieje dokładnie jedno zanurzenie ciał $\iota:\Q\to F$ zachowujące jedynkę. Zachowuje ono i odzwierciedla porządek. W szczególności $F$ ma charakterystykę zero.
\end{twierdzenie}
\begin{proof}
Niech $n_F$ oznacza sumę $n$ kopii $1_F$, przy czym $0_F$ jest zerem ciała. Indukcja i rozdzielność dają $(n+m)_F=n_F+m_F$ i $(nm)_F=n_Fm_F$: dla mnożenia krok indukcji ma postać $n_F(m_F+1)=n_Fm_F+n_F$. Ponieważ $1_F>0$, ciąg tych sum jest ściśle rosnący w sensie porządku indeksów naturalnych, więc $n_F>0$ dla $n\ge1$. W szczególności żadna dodatnia wielokrotność jedynki nie jest zerem.
Dla liczby całkowitej reprezentowanej przez różnicę $a-b$ liczb naturalnych przyjmujemy $(a-b)_F=a_F-b_F$. Jeśli $a-b=c-d$, to $a+d=b+c$, więc $a_F-b_F=c_F-d_F$. Dodawanie i mnożenie zachowują się zgodnie z rachunkiem różnic. Jeżeli $a-b>0$, to $a=b+k$ dla pewnego $k\ge1$, stąd $a_F-b_F=k_F>0$. To daje zachowanie i odzwierciedlanie porządku na $\Z$, w szczególności iniektywność.
Dla $p\in\Z$, $q\in\N_+$ definiujemy $\iota(p/q)=p_F/q_F$. Mianownik jest dodatni, a równość $p/q=r/s$ jest równoważna $ps=rq$; zanurzenie liczb całkowitych zachowuje tę równość, więc definicja nie zależy od reprezentanta. Wzory na sumę i iloczyn ułamków pokazują zachowanie działań. Znak $p_F/q_F$ jest znakiem $p$, dlatego $\iota$ zachowuje i odzwierciedla porządek (stosujemy to do różnicy dwóch ułamków). Każdy homomorfizm zachowujący jedynkę musi zachowywać jej sumy, przeciwności i odwrotności elementów niezerowych, więc musi mieć podany wzór.
\end{proof}
Odtąd w dowolnym ciele uporządkowanym utożsamiamy $\N,\Z,\Q$ z ich obrazami przez to zanurzenie.

\section{Własność Archimedesa}
\begin{definicja}
\label{ma1:def:archimedean}
Ciało uporządkowane $F$ nazywamy archimedesowym, jeżeli dla każdego $x\in F$ istnieje $n\in\N$ takie, że $x<n$.
\end{definicja}
\begin{twierdzenie}[Równoważne postacie własności Archimedesa]
\label{ma1:thm:archimedean-equivalences}
Dla ciała uporządkowanego $F$ następujące warunki są równoważne:
\begin{enumerate}[label=(\roman*)]
\item $F$ jest archimedesowe;
\item dla każdego $\varepsilon>0$ w $F$ istnieje $n\in\N_+$ z $1/n<\varepsilon$;
\item dla każdego $x\in F$ istnieje dokładnie jedno $m\in\Z$ takie, że $m\le x<m+1$;
\item dla wszystkich $x,y\in F$ z $x<y$ istnieje $r\in\Q$ z $x<r<y$.
\end{enumerate}
\end{twierdzenie}
\begin{proof}
(i)$\Rightarrow$(ii): wybieramy $n>1/\varepsilon$ i stosujemy regułę odwracania dodatnich nierówności.
(ii)$\Rightarrow$(i): dla $x>0$ stosujemy (ii) do $1/x$; dla $x\le0$ wystarczy $n=1$.
(i)$\Rightarrow$(iii): wybieramy $N\in\N_+$ z $N>|x|$. Zbiór liczb całkowitych $k$ spełniających $-N\le k\le N$ i $k\le x$ jest niepusty, skończony, ma więc element największy $m$. Ponieważ $N>x$, mamy $m<N$. Jeśli $m+1\le x$, to $m+1$ należałoby do tego samego zbioru, wbrew maksymalności. Jednoznaczność: dla $m<m'$ całkowitych mamy $m+1\le m'$, co wyklucza jednoczesne nierówności $x<m+1$ i $m'\le x$.
(iii)$\Rightarrow$(i): dla $m\le x<m+1$ wybieramy naturalne $n>|m|+1$ (np. $|m|+2$).
(i)$\Rightarrow$(iv): z (ii) wybieramy $n$ z $1/n<y-x$. Z (iii) bierzemy $m\in\Z$ z $m\le nx<m+1$. Wtedy $x<(m+1)/n\le x+1/n<y$.
(iv)$\Rightarrow$(i): między $x$ i $x+1$ istnieje $r=p/q\in\Q$. W $\Q$ istnieje naturalne $n>r$, np. $|p|+1$ dla $q\ge1$, i wtedy $x<r<n$ także w $F$.
\end{proof}
\begin{definicja}[Część całkowita]
Dla $x$ w archimedesowym ciele uporządkowanym oznaczamy przez $\lfloor x\rfloor$ jedyną liczbę całkowitą $m$ z $m\le x<m+1$, a przez $\{x\}=x-\lfloor x\rfloor$ jego część ułamkową.
\end{definicja}
\begin{przyklad}
Dla $x=-7/3$ mamy $\lfloor x\rfloor=-3$ oraz $\{x\}=2/3$. Nie należy mylić części całkowitej z obcięciem ułamka w kierunku zera.
Dla każdego $x$ i $n\in\N_+$ liczba $r_n=\lfloor nx\rfloor/n$ spełnia $r_n\le x<r_n+1/n$. Jest to przybliżenie wymierne z jawnym oszacowaniem błędu, bez użycia teorii granic.
\end{przyklad}
\begin{przyklad}[Ciało niearchimedesowe]
\label{ma1:ex:nonarchimedean}
Niech $F=\Q(t)$ będzie ciałem funkcji wymiernych w formalnej zmiennej $t$. Niezerowy element $p(t)/q(t)$ uznajemy za dodatni, gdy iloraz współczynników wiodących wielomianów $p$ i $q$ jest dodatni.
\end{przyklad}
Sprawdźmy, że określa to stożek dodatni. Jeżeli $p/q=r/s\ne0$, to $ps=rq$; porównanie współczynników wiodących daje tę samą wartość ilorazu, niezależnie od reprezentacji. Każdy niezerowy element ma dokładnie jeden z dwóch znaków, a iloczyn dodatnich ma dodatni iloraz współczynników wiodących. Przy dodawaniu można wybrać mianowniki $q,s$ o dodatnich współczynnikach wiodących. Wówczas również liczniki $p,r$ dodatnich funkcji mają dodatnie współczynniki wiodące. Wielomiany $ps$ i $rq$ mają dodatnie współczynniki wiodące; ich suma nie może utracić najwyższego stopnia przez skracanie, nawet gdy stopnie są równe. Zatem $(ps+rq)/(qs)>0$. Twierdzenie \ref{ma1:thm:positive-cone} daje porządek ciała.
Dla każdego $n\in\N$ wielomian $t-n$ ma współczynnik wiodący $1$, więc $t>n$. Ciało nie jest archimedesowe. Ponadto $0<1/t<1/n$ dla wszystkich $n\ge1$. Jest to przykład dodatniego elementu mniejszego od każdej dodatniej liczby wymiernej: dla $p/q>0$ mamy $p\ge1$ i $1/t<1/q\le p/q$. Nie utożsamiamy tego ciała z ciałem liczb hiperrealnych.

\section{Zupełność porządkowa}
\begin{definicja}[Kresy i zupełność]
\label{ma1:def:order-complete}
Niech $F$ będzie ciałem uporządkowanym, a $E\subseteq F$ niepustym zbiorem. Element $u\in F$ jest ograniczeniem górnym $E$, jeśli $x\le u$ dla każdego $x\in E$. Kres górny $\sup E$, jeśli istnieje, jest najmniejszym z ograniczeń górnych. Analogicznie ograniczenie dolne $l$ spełnia $l\le x$ dla każdego $x\in E$, a kres dolny $\inf E$ jest największym z ograniczeń dolnych. Ciało $F$ nazywamy zupełnym porządkowo, jeśli każdy niepusty podzbiór ograniczony z góry posiada kres górny w $F$.
\end{definicja}
\begin{lemat}[Kryterium kresu]
\label{ma1:lem:sup-criterion}
Niech $F$ będzie ciałem uporządkowanym, $E\subseteq F$ niepustym zbiorem i $s\in F$. Wówczas $s=\sup E$ wtedy i tylko wtedy, gdy $s$ jest ograniczeniem górnym $E$ oraz dla każdego $\varepsilon>0$ w $F$ istnieje $x\in E$ z $s-\varepsilon<x$. Analogicznie $s=\inf E$ wtedy i tylko wtedy, gdy $s$ jest ograniczeniem dolnym i dla każdego $\varepsilon>0$ istnieje $x\in E$ z $x<s+\varepsilon$.
\end{lemat}
\begin{proof}
Jeśli $s=\sup E$, to $s-\varepsilon<s$ nie jest ograniczeniem górnym, więc istnieje żądany $x$. Jeśli warunki zachodzą, dowolne $u<s$ nie jest ograniczeniem górnym, bo dla $\varepsilon=s-u>0$ istnieje $x>u$ w $E$. Zatem $s$ jest najmniejszym ograniczeniem górnym. Dla kresu dolnego stosujemy ten wynik do $-E=\{-x:x\in E\}$: negacja odwraca porządek, więc $s=\inf E$ dokładnie wtedy, gdy $-s=\sup(-E)$.
\end{proof}
\begin{wniosek}
\label{ma1:cor:inf-exists}
Ciało uporządkowane $F$ jest zupełne porządkowo wtedy i tylko wtedy, gdy każdy jego niepusty podzbiór ograniczony z dołu ma kres dolny.
\end{wniosek}
\begin{proof}
Jeśli $E$ jest ograniczony z dołu, $-E$ jest ograniczony z góry i $\inf E=-\sup(-E)$. Odwrotny kierunek wynika analogicznie z $\sup E=-\inf(-E)$.
\end{proof}
\begin{twierdzenie}[Działania na kresach]
\label{ma1:thm:sup-arithmetic}
Niech $E,G$ będą niepustymi podzbiorami ciała uporządkowanego $F$, mającymi kresy górne $s=\sup E$, $t=\sup G$. Wtedy $\sup(E+G)=s+t$. Dla $\lambda>0$ zachodzi $\sup(\lambda E)=\lambda s$, a dla $\lambda=0$ kres wynosi zero. Jeżeli $E,G\subseteq\{x:x\ge0\}$, to $\sup(EG)=st$. Jeśli $E$ ma kres dolny $l$ i $\lambda<0$, to $\sup(\lambda E)=\lambda l$. Zbiory $E+G,EG,\lambda E$ oznaczają odpowiednio zbiory wszystkich sum, iloczynów i wielokrotności elementów.
\end{twierdzenie}
\begin{proof}
Dla sumy $s+t$ jest ograniczeniem górnym. Dla $\varepsilon>0$ wybieramy $e>s-\varepsilon/2$ w $E$ i $g>t-\varepsilon/2$ w $G$. Wtedy $e+g>s+t-\varepsilon$ i stosujemy lemat \ref{ma1:lem:sup-criterion}. Dla $\lambda>0$ ograniczenie górne $\lambda s$ jest oczywiste, a wybór $e>s-\varepsilon/\lambda$ daje $\lambda e>\lambda s-\varepsilon$. Dla zera zbiór wielokrotności jest równy $\{0\}$. Dla $\lambda<0$ negacja zamienia kres dolny na górny, a następnie stosujemy wynik dla dodatniego $-\lambda$ do $-E$.
Dla iloczynu nieujemnych zbiorów $eg\le st$. Niech $M=1+s+t>0$. Dla $\varepsilon>0$ wybieramy $e>s-\varepsilon/M$ i $g>t-\varepsilon/M$. Wówczas $0\le s-e<\varepsilon/M$ i $0\le t-g<\varepsilon/M$, a
\[
 0\le st-eg=t(s-e)+e(t-g)
 \le t(s-e)+s(t-g)<\varepsilon
\]
(gdy $s+t=0$, różnica jest zerem; w pozostałych przypadkach ostatnie oszacowanie wynika z $(s+t)/M<1$). Kryterium kresu kończy dowód.
\end{proof}
\begin{twierdzenie}[Zupełność implikuje własność Archimedesa]
\label{ma1:thm:complete-archimedean}
Każde zupełne porządkowo ciało uporządkowane jest archimedesowe.
\end{twierdzenie}
\begin{proof}
Gdyby $\N$ była ograniczona z góry w $F$, miałaby kres górny $s$. Ponieważ $s-1<s$, istnieje $n\in\N$ z $n>s-1$. Wówczas $n+1\in\N$ oraz $n+1>s$, sprzecznie z tym, że $s$ jest ograniczeniem górnym.
\end{proof}
\begin{przyklad}
$\Q$ jest archimedesowe, lecz nie jest zupełne porządkowo: brak kresu górnego zbioru nieujemnych liczb wymiernych o kwadracie mniejszym od $2$ wykazaliśmy w rozdziale \ref{ma1:ch:number-systems}. Z kolei ciało z przykładu \ref{ma1:ex:nonarchimedean} nie jest zupełne, na mocy twierdzenia \ref{ma1:thm:complete-archimedean}. Ciało $\R$ z rozdziału \ref{ma1:ch:real-construction} jest zupełne, a więc archimedesowe.
\end{przyklad}

\section{Konsekwencje zupełności liczb rzeczywistych}
\begin{twierdzenie}[Przedziały zstępujące]
\label{ma1:thm:nested-intervals}
Niech $(I_n)_{n\ge1}$ będzie rodziną przedziałów $I_n=[a_n,b_n]\subseteq\R$, gdzie $a_n\le b_n$ i $I_{n+1}\subseteq I_n$ dla każdego $n\ge1$. Wtedy $\bigcap_{n\ge1}I_n\ne\varnothing$. Jeżeli dodatkowo dla każdego $\varepsilon>0$ istnieje $n\ge1$ z $b_n-a_n<\varepsilon$, przecięcie zawiera dokładnie jeden punkt.
\end{twierdzenie}
\begin{proof}
Z inkluzji wynika $a_n\le a_{n+1}\le b_{n+1}\le b_n$. Dla dowolnych $k,n$ i $m=\max(k,n)$ mamy $a_k\le a_m\le b_m\le b_n$. W szczególności $b_1$ ogranicza z góry wszystkie $a_k$. Niech $s=\sup\{a_k:k\ge1\}$. Dla każdego $n$ mamy $a_n\le s\le b_n$, ponieważ $b_n$ również jest ograniczeniem górnym zbioru lewych końców. Stąd $s\in I_n$ dla wszystkich $n$.
Gdyby przecięcie zawierało dwa różne punkty $x<y$, to $y-x\le b_n-a_n$ dla każdego $n$, co przeczy założeniu zastosowanemu do $\varepsilon=y-x>0$.
\end{proof}
\begin{przyklad}
Dla $I_n=[-1/n,1/n]$ przecięcie jest równe $\{0\}$. Natomiast przedziały otwarte $(0,1/n)$ tworzą rodzinę zstępującą o pustym przecięciu: dla każdego $x>0$ istnieje $n$ z $1/n<x$. Bez warunku domknięcia twierdzenie nie zachodzi. Przedziały $[0,1+1/n]$ pokazują, że bez warunku dowolnie małej długości przecięcie może być większe od jednego punktu (jest równe $[0,1]$).
\end{przyklad}

\begin{figure}[htbp]
\centering
\begin{tikzpicture}[x=2.0cm,y=0.65cm,>=Stealth]
\draw[->] (-1.7,0)--(1.7,0) node[right] {$x$};
\draw (0,-.12)--(0,.12) node[below=3pt] {$0$};
\foreach \k/\w in {1/1,2/0.5,3/0.3333}{
  \draw[blue!65!black,line width=1.1pt] (-\w,\k)--(\w,\k);
  \fill[blue!65!black] (-\w,\k) circle (1.5pt) (\w,\k) circle (1.5pt);
  \node[anchor=west] at (1.15,\k) {$I_{\k}$};
}
\draw[densely dashed,gray] (0,.2)--(0,3.35);
\end{tikzpicture}
\caption{Przedziały $I_n=[-1/n,1/n]$ mają wspólny punkt $0$, a ich długości są dowolnie małe.}
\label{ma1:fig:nested-intervals}
\end{figure}

\begin{lemat}[Oszacowanie przyrostu potęgi]
\label{ma1:lem:power-bound}
Niech $F$ będzie ciałem uporządkowanym, $n\in\N_+$, $M>0$ oraz $0\le x\le y\le M$. Wówczas
\[
 0\le y^n-x^n=(y-x)\sum_{k=0}^{n-1}y^{n-1-k}x^k
 \le nM^{n-1}(y-x).
\]
Jeżeli $x<y$, to $x^n<y^n$.
\end{lemat}
\begin{proof}
Równość wynika z wymnożenia sumy i skrócenia kolejnych wyrazów. Każdy składnik sumy jest nieujemny i nie większy od $M^{n-1}$. Jeśli $x<y$, to $y>0$, więc pierwszy składnik $y^{n-1}$ jest dodatni (dla $n=1$ jest równy $1$). Iloczyn przez $y-x>0$ jest dodatni.
\end{proof}
\begin{twierdzenie}[Istnienie i jednoznaczność pierwiastka]
\label{ma1:thm:positive-root}
Niech $a\ge0$ będzie liczbą rzeczywistą, a $n\in\N_+$. Istnieje dokładnie jedno $s\ge0$ takie, że $s^n=a$.
\end{twierdzenie}
\begin{proof}
Dla $a=0$ wystarczy $s=0$, a jednoznaczność wynika z dodatniości potęg liczb dodatnich. Niech $a>0$. Zbiór $E=\{x\in\R:x\ge0,\ x^n\le a\}$ zawiera zero i jest ograniczony z góry przez $M=1+a$: jeśli $x>M$, to $x>1$ i $x^n\ge x>a$. Istnieje więc $s=\sup E$, przy czym $0\le s\le M$. Ponadto $s>0$, ponieważ $e=\min(1,a)/2>0$ spełnia $e^n\le e\le a$.
Przyjmijmy $C=n(M+1)^{n-1}>0$. Jeśli $s^n<a$, wybieramy $0<h<\min(1,(a-s^n)/C)$, np. połowę tej dodatniej liczby. Z lematu \ref{ma1:lem:power-bound} dla $s+h\le M+1$ wynika
$(s+h)^n-s^n\le Ch<a-s^n$, a więc $s+h\in E$, sprzecznie z ograniczeniem górnym $s$.
Jeśli $s^n>a$, wybieramy $0<h<\min(s,(s^n-a)/C)$. Wtedy $s^n-(s-h)^n\le Ch<s^n-a$, zatem $(s-h)^n>a$. Ścisła monotoniczność potęgi na zbiorze nieujemnych liczb oznacza, że żaden $x\in E$ nie może spełniać $x\ge s-h$. Liczba $s-h$ jest więc ograniczeniem górnym $E$ mniejszym od $s$, sprzeczność. Pozostaje $s^n=a$. Jednoznaczność ponownie wynika ze ścisłej monotoniczności.
\end{proof}
Oznaczamy tę liczbę przez $\sqrt[n]{a}$, a dla $n=2$ przez $\sqrt a$. Dla $q\in\Q$, $q\ge0$, ścisła monotoniczność kwadratu daje $q<\sqrt2$ dokładnie wtedy, gdy $q^2<2$. Wszystkie ujemne $q$ także leżą poniżej $\sqrt2$. To identyfikuje przekrój z końca rozdziału \ref{ma1:ch:real-construction} jako liczbę $\sqrt2$.
\begin{wniosek}[Gęstość liczb niewymiernych]
\label{ma1:cor:irrational-density}
Dla dowolnych liczb rzeczywistych $x<y$ istnieje liczba niewymierna $z$ spełniająca $x<z<y$.
\end{wniosek}
\begin{proof}
Liczba $\sqrt2$ nie jest wymierna, gdyż w rozdziale \ref{ma1:ch:number-systems} wykazaliśmy, że żadna liczba wymierna nie ma kwadratu równego $2$. Z gęstości $\Q$ wybieramy $r\in\Q$ takie, że $x-\sqrt2<r<y-\sqrt2$. Kładziemy $z=r+\sqrt2$. Jeśli $z$ byłaby wymierna, to $\sqrt2=z-r$ także byłaby wymierna, sprzeczność.
\end{proof}

\section{Jednoznaczność zupełnego ciała uporządkowanego}
\begin{twierdzenie}
\label{ma1:thm:real-uniqueness}
Dla każdego zupełnego porządkowo ciała uporządkowanego $F$ istnieje dokładnie jeden izomorfizm ciał uporządkowanych $J:F\to\R$. Po utożsamieniu $\R$ z konstrukcją Dedekinda i $\Q$ z kanonicznym podciałem $F$ ma on postać $J(x)=\{q\in\Q:q<x\}$.
\end{twierdzenie}
\begin{proof}
\emph{Przekroje i porządek.} Ciało $F$ jest archimedesowe z twierdzenia \ref{ma1:thm:complete-archimedean}. Dla każdego $x\in F$ istnieje $n\in\N_+$ z $n>|x|$. Zbiór $J(x)$ zawiera $-n$, nie zawiera $n$, jest domknięty w dół, a dla $q<x$ gęstość $\Q$ daje $q<r<x$ z $r\in\Q$. Zatem $J(x)\in\R_D$. Dla $x<y$ wybieramy $q\in\Q$ z $x<q<y$; wtedy $J(x)\subsetneq J(y)$, bo $q\in J(y)\setminus J(x)$. To daje zachowanie i odzwierciedlanie porządku.
Dla przekroju $A\in\R_D$ traktujemy $A$ jako podzbiór $F$ przez kanoniczne zanurzenie $\Q$. Jest on niepusty i ograniczony z góry przez każdy $u\in\Q\setminus A$. Niech $x=\sup_F A$. Dla $q\in A$ istnieje $q'>q$ w $A$, więc $q<x$. Odwrotnie, jeśli $q\in\Q$ i $q<x$, to $q$ nie jest ograniczeniem górnym $A$, więc istnieje $a\in A$ z $q<a$ i stąd $q\in A$. Otrzymujemy $J(x)=A$; $J$ jest bijekcją.
\emph{Dodawanie.} Jeśli $p<x$, $q<y$ są wymierne, to $p+q<x+y$, zatem $J(x)+J(y)\subseteq J(x+y)$. Dla wymiernego $r<x+y$ wybieramy z gęstości $\Q$ liczbę $p$ z $r-y<p<x$. Wtedy $r-p<y$ i $r=p+(r-p)\in J(x)+J(y)$. Zatem $J(x+y)=J(x)+J(y)$. W szczególności $J(0)=D_0$ i $J(-x)=-J(x)$; również $J(1)=D_1$.
\emph{Mnożenie.} Dla $x,y>0$ z $p<x$, $q<y$, $p,q>0$ wynika $pq<xy$. Z definicji dodatniego iloczynu przekrojów $J(x)\odot J(y)\subseteq J(xy)$. Dla $0\le r<xy$ wymiernego wybieramy wymierne $p$ z $r/y<p<x$ i następnie $q$ z $r/p<q<y$. Obie liczby są dodatnie, $r<pq$ i zatem $r\in J(x)\odot J(y)$. Liczby ujemne należą do obu przekrojów. To dowodzi równości. Gdy jeden czynnik jest zerem, obie strony są zerem; dla ujemnych czynników stosujemy zachowanie negacji i prawa znaków w obu ciałach. Otrzymaliśmy izomorfizm ciał uporządkowanych.
\emph{Jednoznaczność.} Każdy izomorfizm ciał uporządkowanych $K:F\to\R_D$ zachowuje jedynkę i kanoniczne liczby wymierne z twierdzenia \ref{ma1:thm:rational-embedding}. Dla dowolnego $x\in F$ oraz $q\in\Q$ mamy $q<x$ wtedy i tylko wtedy, gdy $D_q<K(x)$. Dla dowolnego przekroju $B$ równoważność $D_q<B\Leftrightarrow q\in B$ wynika z definicji: jeśli $q\in B$, to $D_q\subsetneq B$; jeśli $D_q\subsetneq B$, istnieje $r\in B\setminus D_q$, a więc $r\ge q$ i $q\in B$. Zatem $K(x)$ i $J(x)$ zawierają dokładnie te same liczby wymierne, czyli są równe.
\end{proof}
\begin{uwaga}
Zupełność porządkowa nie oznacza, że dowolny zbiór ma kres górny: wymagamy niepustości i ograniczoności z góry. W następnym rozdziale rozważymy rozszerzenia systemu liczb rzeczywistych, a później zbieżność w przestrzeniach metrycznych.
\end{uwaga}
